\section{Model, notation and the index targets in use}
\label{sec:model}

\subsection{Data model}

The data are maps $d_b$ in bands $b = 1,\dots,n_{\rm b}$, each a vector over elements $p$ (pixels,
and Stokes parameters in polarization), in thermodynamic units. The parametric model is
\begin{equation}
  d_b = \sum_c m_{b,c}(\beta)\, a_c + n_b , \qquad n_b \sim \mathcal{N}(0, \mathsf{N}_b),
  \label{eq:model}
\end{equation}
with one amplitude $a_c(p)$ per component $c$ and element, a mixing coefficient $m_{b,c}$ that is the
component spectrum averaged over the bandpass of band $b$, and spectral parameters $\beta$ that are
shared by many elements: one value for the sky, one per region, or a field at lower resolution than
the amplitudes. Stacking bands and components, $d = \M(\beta)\,a + n$ with $a$ of length $k$. We write
\begin{equation}
  \A = \M\T \Ni \M , \qquad \hat a = \A^{-1} \M\T \Ni d , \qquad
  \Pm = \Ni - \Ni \M \A^{-1} \M\T \Ni ,
\end{equation}
all functions of $\beta$, so that
\begin{equation}
  \chi^2(\beta, a) = \chimin(\beta) + (a - \hat a)\T \A\, (a - \hat a), \qquad \chimin(\beta) = d\T \Pm d .
  \label{eq:chisplit}
\end{equation}
For the power-law component that carries the index, $m_{b}(\beta)$ is the band average of
$(\nu/\nuref)^{\beta}$ in brightness temperature, converted to thermodynamic units, so that the
amplitude is the brightness temperature at the reference frequency $\nuref$. The reference frequency
is a convention: it changes the meaning of $a$ and not the model.

The amplitudes are what statisticians call incidental parameters \citep{neyman1948}: their number
grows with the data, and the parameters of interest, the indices, are few. Everything in this paper
is a consequence of how the incidental parameters are removed.

\subsection{Four ways to remove the amplitudes}

\paragraph{Joint posterior with flat amplitude priors.} With $p(a) \propto 1$ the joint posterior is
$p(\beta)\exp[-\chi^2(\beta,a)/2]$ and integrating Eq.~(\ref{eq:chisplit}) over $a$ gives the index
marginal
\begin{equation}
  \mathcal{L}_{\rm flat}(\beta) \propto \exp\!\left[-\tfrac12 \chimin(\beta)\right] \det \A(\beta)^{-1/2} .
  \label{eq:Lflat}
\end{equation}
A Gibbs sampler that draws $a \mid \beta$ with its fluctuation and then moves $\beta$ on $\chi^2$ at
fixed $a$ has exactly this marginal for $\beta$, whether or not it mixes well enough to show it
(checked numerically and re-derived independently, \exref{E012}). The determinant is the volume
of amplitude space allowed at a given $\beta$. It is an implicit prior on the index and the subject of
Sect.~\ref{sec:flat}.

\paragraph{Profile likelihood.} Maximizing over the amplitudes instead of integrating gives
\begin{equation}
  \mathcal{L}_{\rm prof}(\beta) \propto \exp\!\left[-\tfrac12 \chimin(\beta)\right],
  \label{eq:Lprof}
\end{equation}
the spectral likelihood of \citet{stompor2009} (their Eq.~9; their Eq.~10 is the marginal form).
As a Bayesian target it is the marginal under the amplitude prior
$p(a \mid \beta) \propto \det \A(\beta)^{1/2}$, a Jeffreys prior for the amplitude block conditional on
the index. It is not the marginal under the joint Jeffreys prior of $(a, \beta)$, which depends on $a$.

\paragraph{Gaussian amplitude prior.} With $a \sim \mathcal{N}(\mu, \mathsf{S})$ the marginal keeps the
form of Eq.~(\ref{eq:Lflat}) with $\A \to \A + \mathsf{S}^{-1}$ and the prior mean entering $\chimin$.
We use two cases. A zero-mean prior with one free scale per index region,
$a_{\rm s}(p) \sim \mathcal{N}(0, \tau_r^2)$ for $p$ in region $r$, with $\tau_r$ maximized or
marginalized under a prior flat in $\ln\tau_r$. And, for intensity, a prior whose mean and scale are
tied to a template, with hyperparameters inferred. Under a change of reference frequency $\tau_r$
rescales by $(\nuref'/\nuref)^{\beta_r}$, and both maximization and a log-flat prior are invariant
under that rescaling, so the first case does not depend on the reference frequency.

\paragraph{Re-solving the amplitudes inside the index move.} A fourth route is procedural: propose
$\beta'$, solve for new amplitudes at $\beta'$, and accept on the $\chi^2$ of the new pair. What this
samples depends on how the new amplitudes are produced, and is the subject of
Sect.~\ref{sec:legacy}.

\subsection{Information and score}

Let $\M_r = \partial \M/\partial\beta_r$ for index parameter $r$. For data that follow
Eq.~(\ref{eq:model}) with true amplitudes $a$, the profile score
$s_r = -\tfrac12\,\partial\chimin/\partial\beta_r$ has zero mean and covariance
\begin{equation}
  \mathrm{Cov}(s_r, s_{r'}) = \J_{rr'} + \Tm_{rr'}, \qquad
  \J_{rr'} = (\M_r a)\T \Pm\, (\M_{r'} a), \qquad
  \Tm_{rr'} = \tr\!\left(\A^{-1} \M_r\T \Pm\, \M_{r'}\right),
  \label{eq:JT}
\end{equation}
while the expected curvature of $\tfrac12\chimin$ is $\J$ alone (derivation in
App.~\ref{app:sandwich}). $\J$ is the Fisher information on the indices with the amplitudes
profiled out. $\Tm$ does not depend on the signal: it is the part of the score variance that the
noise in the fitted amplitudes contributes. For one power-law component, diagonal noise with weights
$w_{pb} = 1/\sigma_{pb}^2$ and one region, these reduce to
\begin{equation}
  \J = \sum_p u_p V_p, \qquad \Tm = \sum_p V_p, \qquad
  u_p = a_p^2 \sum_b w_{pb} m_b^2, \qquad
  V_p = \mathrm{Var}_{w m^2}(\ln \nu)_p ,
  \label{eq:uV}
\end{equation}
with $u_p$ the squared signal-to-noise ratio of amplitude $p$ and $V_p$ the variance of $\ln\nu$ over
the bands under the weights $w_{pb} m_b^2$: the lever arm.
